% ================================================================
% Chapter 15: Repertoire Curation and Museum Theory
% ================================================================

\chapter{Repertoire Curation and Museum Theory}
\label{ch:curation}

\epigraph{The cultural arbitrary which a system of education transmits,
reproduces, and inculcates is \emph{always} the cultural arbitrary of the
dominant classes.}{Pierre Bourdieu and Jean-Claude Passeron,
\textit{La Reproduction} (1970)}

% ----------------------------------------------------------------
\section{Introduction}
% ----------------------------------------------------------------

The preceding chapters have treated the receiver as a fixed entity---a
repertoire passively awaiting signals.  But receivers are not passive.
Every mind \emph{curates} its repertoire: selecting which ideas to retain,
which to discard, and which new ideas to admit.  The museum curator
choosing which artifacts to display, the teacher designing a syllabus, the
censor deciding which books to ban, the autodidact selecting their next
reading---all are engaged in the same structural operation: filtering and
expanding a repertoire according to some criterion.

Bourdieu's concept of \emph{cultural capital}---the accumulated, curated
stock of ideas that determines one's interpretive capacity---is the
sociological correlate of the repertoire.  Cultural capital is not merely
\emph{possessed}; it is actively \emph{managed}.  The bourgeois family
curates its children's education; the academy curates its canon; the state
curates the national curriculum.  Each act of curation is simultaneously
an act of inclusion and exclusion, of memory and forgetting.\footnote{On the
politics of canon formation, see John Guillory, \textit{Cultural Capital:
The Problem of Literary Canon Formation} (1993).  Guillory argues that the
literary canon is a form of cultural capital whose curation serves class
reproduction---a thesis that our formal framework substantiates.}

This chapter formalizes repertoire curation and expansion.  We introduce
$\operatorname{curate}$---filtering a repertoire by a predicate---and
$\operatorname{expand}$---adding a new idea to the front.  We then prove
eighteen theorems that connect these operations to museum theory, education
policy, censorship dynamics, and the political economy of knowledge.

\subsection*{The Museum Without Walls}

Andr\'{e} Malraux's \textit{Le Mus\'{e}e imaginaire} (1947) introduced the
revolutionary idea of a ``museum without walls'': through photographic
reproduction, every artwork becomes available to every viewer regardless
of geography.  Malraux argued that this technological shift fundamentally
altered the nature of curation---the curator was no longer the gatekeeper
of a physical collection but an \emph{editor} assembling a virtual
repertoire from the totality of world art.  In our formal framework,
Malraux's insight translates directly: the advent of cheap reproduction
means that the set $\IS$ available for curation expands dramatically, but
the curatorial predicate $p$ remains the structural bottleneck.  The
museum without walls does not eliminate curation; it makes curation
\emph{more necessary}, because the uncurated set has grown beyond any
individual's capacity to survey.\footnote{Malraux's argument anticipates
the information-overload problems of the digital age.  When every image
is available online, the curatorial function becomes \emph{more} valuable,
not less.  The ``museum without walls'' is, in our terms, a repertoire
of effectively infinite cardinality---making the choice of predicate $p$
the decisive act.}

Tony Bennett's \textit{The Birth of the Museum} (1995) offered a
complementary perspective, analyzing museums as \emph{disciplinary
institutions} in the Foucauldian sense.  For Bennett, the museum does not
merely curate knowledge; it produces docile subjects by training visitors
to see the world through the museum's organizing categories.  The
repertoire $\mathbf{R}$ displayed by a museum is not a neutral selection
but a \emph{technology of citizenship}, shaping the visitor's own
repertoire through repeated exposure.  Bennett's analysis maps onto our
formal framework as follows: the museum's displayed repertoire
$\operatorname{curate}(\mathbf{R}_{\text{total}}, p_{\text{institutional}})$
becomes the basis for the visitor's expansion operations.  The visitor
expands their personal repertoire $\mathbf{R}_{\text{visitor}}$ by
elements drawn from the museum's curated display, so that institutional
curation upstream constrains individual expansion downstream.

\subsection*{Education as Repertoire Building}

The connection between curation and education runs deeper than analogy.
John Dewey's progressive pedagogy, articulated in \textit{Democracy and
Education} (1916), conceived of education as the guided expansion of
experience---a process we formalize as iterated application of
$\operatorname{expand}$.  Dewey rejected the ``banking model'' in which
knowledge is deposited into passive minds; instead, he argued that
education must be an active process of inquiry in which the learner
selects (curates) their experiences.  The Deweyan student is simultaneously
expanding and curating: each new experience is an expansion, and the
reflective evaluation of that experience is a curation step that
determines which elements persist in the repertoire.

Paulo Freire radicalized Dewey's insight in \textit{Pedagogy of the
Oppressed} (1968), arguing that the ``banking model'' of education is
itself a form of oppressive curation.  When the teacher controls all
expansion operations---deciding which ideas enter the student's
repertoire---and the student has no voice in the curatorial predicate $p$,
education becomes domination.  Freire's alternative, ``problem-posing
education,'' gives the student co-authorship over both operations:
the student participates in deciding what to learn
($\operatorname{expand}$) and how to evaluate what has been learned
($\operatorname{curate}$).  In our formal terms, Freirean liberation
consists in democratizing the predicate $p$---transferring curatorial
authority from the oppressor to the community of learners.

\subsection*{Censorship as Negative Curation}

If education is positive curation---selecting \emph{for} inclusion---then
censorship is negative curation: selecting \emph{for exclusion}.  The
\textit{Index Librorum Prohibitorum}, maintained by the Catholic Church
from 1559 to 1966, was among the most systematic curatorial operations
in history.  The Index defined a predicate
$p_{\text{Index}}(r) = \texttt{false}$ for thousands of works by
Copernicus, Galileo, Descartes, Hume, Voltaire, and others, effectively
curating the Catholic intellectual repertoire by negation.  The structural
consequence, formalized in our Theorem~\ref{thm:curate-length}, is
clear: $|\operatorname{curate}(\mathbf{R}_{\text{Catholic}},
p_{\text{Index}})| < |\mathbf{R}_{\text{Catholic}}|$.  The censored
repertoire is strictly smaller---and therefore, by
Theorem~\ref{thm:curated-interp}, generates strictly fewer
interpretations of any signal.

Contemporary book banning---from school library challenges in the United
States to internet censorship in authoritarian regimes---follows the same
structural logic.  Each banned book is an element removed from the
communal repertoire; each act of censorship narrows the society's
interpretive capacity.  The formal framework makes precise what civil
libertarians assert intuitively: censorship makes us collectively dumber,
in the strict sense of reducing the number of available interpretations.

\subsection*{The Canon Wars}

The ``canon wars'' of the 1980s and 1990s---the debate over whether
Western universities should center their curricula on the ``Great Books''
of the Western tradition or expand to include non-Western and minority
voices---are, in IDT terms, a dispute about the appropriate predicate
$p$.  Defenders of the traditional canon (Allan Bloom's \textit{The
Closing of the American Mind}, 1987; Harold Bloom's \textit{The Western
Canon}, 1994) argued for a predicate that privileged aesthetic excellence
as judged by established critical traditions.  Advocates of multicultural
expansion (Henry Louis Gates Jr., bell hooks, Edward Said) argued that
the traditional predicate systematically excluded the cultural
productions of women, people of color, and non-Western societies.

Our formal framework illuminates both sides of the debate.  The
traditionalists are correct that some curatorial predicate is necessary
(Theorem~\ref{thm:universal-curate}: universal acceptance is not
curation at all, and the resulting repertoire is unnavigable).  The
multiculturalists are correct that a narrow predicate reduces
interpretive capacity (Theorem~\ref{thm:curated-interp}: fewer retained
ideas means fewer interpretations of any signal).  The resolution, if
there is one, lies in designing predicates that balance navigability
against breadth---a problem our theorems frame precisely but do not
solve, because the choice of $p$ is irreducibly political.

\subsection*{Digital Curation and the Paradox of Choice}

The digital age has introduced a new species of curation: algorithmic
recommendation.  Spotify's Discover Weekly, YouTube's recommendation
engine, and TikTok's For You page all implement curatorial predicates
$p_{\text{algorithm}}$ that filter the vast ideatic space of available
content into a personalized repertoire for each user.  These algorithms
curate at a scale and speed impossible for any human curator, processing
millions of items to produce a feed of dozens.  The structural properties
are identical to those of the museum curator: the algorithm filters
($\operatorname{curate}$), the user absorbs ($\operatorname{expand}$),
and the resulting repertoire determines interpretive capacity.

Barry Schwartz's \textit{The Paradox of Choice} (2004) identified the
psychological cost of excessive options: when the uncurated set $\IS$ is
too large, individuals become paralyzed, anxious, and dissatisfied.
Schwartz's paradox is, in our terms, a consequence of the gap between
$|\IS|$ and human cognitive capacity to evaluate curatorial predicates.
The algorithm resolves this paradox by imposing a predicate on the user's
behalf---but at the cost of surrendering curatorial autonomy to an opaque
system whose predicate $p_{\text{algorithm}}$ is optimized for engagement
metrics rather than interpretive depth.  The formal question becomes:
does algorithmic curation increase or decrease
$\operatorname{depthSum}(\mathbf{R}_{\text{user}})$?  Our theorems
(Theorem~\ref{thm:curated-depthsum}) show that curation can only
\emph{reduce} total depth; whether the algorithm's curation reduces
depth less than the user's own imperfect curation is an empirical
question that the formal framework poses but cannot answer
\textit{a priori}.

% ----------------------------------------------------------------
\section{Core Definitions}
\label{sec:ch15-defs}
% ----------------------------------------------------------------

\begin{definition}[Curation]
\label{def:curate}
To \emph{curate} a repertoire $\mathbf{R}$ is to filter it by a Boolean
predicate $p$:
\[
  \operatorname{curate}(\mathbf{R}, p)
  \;=\; [r \in \mathbf{R} \mid p(r) = \texttt{true}].
\]
This models selective retention: the museum curator choosing artifacts, the
teacher choosing concepts, the censor choosing permitted ideas.
\end{definition}

\begin{definition}[Expansion]
\label{def:expand}
To \emph{expand} a repertoire $\mathbf{R}$ with a new idea $a$ is to
prepend it:
\[
  \operatorname{expand}(\mathbf{R}, a) \;=\; a :: \mathbf{R}.
\]
This models learning: acquiring a new concept, encountering a new artifact,
absorbing a new cultural input.
\end{definition}

\begin{lstlisting}[caption={Curation and expansion in Lean~4}]
def curate (rep : List I) (p : I -> Bool) : List I :=
  rep.filter p

def expand (rep : List I) (new_idea : I) : List I :=
  new_idea :: rep
\end{lstlisting}

\begin{definition}[Interpretation]
\label{def:interpret-ch15}
When a receiver with repertoire $\mathbf{R}$ encounters signal $s$, they
produce interpretations:
\[
  \operatorname{interpret}(\mathbf{R}, s)
  \;=\; [r_1 \compose s,\; r_2 \compose s,\; \dots,\; r_n \compose s].
\]
\end{definition}

The interplay of curation, expansion, and interpretation is the central
drama of this chapter.  Curation \emph{reduces} the repertoire (and hence
the number and depth of possible interpretations); expansion
\emph{increases} it.  The structural question is: how do these operations
compose, and what invariants do they preserve?

\begin{example}[The Alexandrian Library as Curation Catastrophe]
\label{ex:alexandria}
The Library of Alexandria, founded in the third century BCE under the
Ptolemaic dynasty, represented the ancient world's most ambitious attempt
at \emph{anti-curation}: the aspiration to collect every written work in
existence.  In our formal terms, the Alexandrian project sought to
approximate $\mathbf{R}_{\text{Alexandria}} \approx \IS_{\text{written}}$
---a repertoire equal to the entire space of written ideas.  By
Theorem~\ref{thm:universal-curate}, this amounts to the identity
curation $\operatorname{curate}(\IS, \lambda x.\, \texttt{true}) = \IS$,
which is structurally equivalent to no curation at all.

The destruction of the Library---whether by fire, neglect, or
conquest---constitutes the most famous curation catastrophe in history.
The curatorial predicate shifted from $\lambda x.\, \texttt{true}$
(preserve everything) to something far more restrictive, dictated by the
contingencies of survival: which scrolls happened to be copied, which
survived the fires, which were translated into Arabic or Latin.  The
resulting repertoire $\mathbf{R}_{\text{surviving}}$ is a strict subset
of $\mathbf{R}_{\text{Alexandria}}$, and by
Theorem~\ref{thm:curated-depthsum},
$\operatorname{depthSum}(\mathbf{R}_{\text{surviving}}) \le
\operatorname{depthSum}(\mathbf{R}_{\text{Alexandria}})$.  The lost
depth---the works of Sappho, Aristotle's second book of the
\textit{Poetics}, countless treatises on mathematics, astronomy, and
philosophy---represents an irrecoverable reduction in humanity's
interpretive capacity.
\end{example}

\begin{example}[The Chinese Imperial Examination as Repertoire Filter]
\label{ex:keju}
The Chinese imperial examination system (\textit{k\={e}j\v{u}}),
operational from 605 CE to 1905, was among the longest-running and most
consequential curatorial operations in human history.  The examinations
required mastery of a fixed canon: the Four Books and Five Classics of
Confucian literature.  In formal terms, the examination defined a
curatorial predicate $p_{\text{k\={e}j\v{u}}}$ that accepted ideas
consonant with the Confucian classics and rejected all others.  Every
aspiring official curated their personal repertoire according to this
predicate: $\mathbf{R}_{\text{scholar}} = \operatorname{curate}(
\mathbf{R}_{\text{available}}, p_{\text{k\={e}j\v{u}}})$.

The system's strength was its meritocratic inclusiveness: any man
(and, in rare cases, any woman) could in principle expand their
repertoire sufficiently to pass.  Its weakness was the narrowness of the
predicate: by the late imperial period, the ``eight-legged essay''
format had ossified the curatorial criterion to the point where
$p_{\text{k\={e}j\v{u}}}$ rejected virtually all modern scientific and
technological ideas.  When Western ideas arrived in force in the
nineteenth century, the Qing dynasty possessed a governing class whose
repertoires had been curated for literary elegance and moral philosophy
but systematically stripped of empirical science.  The result was a
catastrophic mismatch between the curated repertoire and the signals
arriving from the industrializing West.
\end{example}

% ----------------------------------------------------------------
\section{Curation Fundamentals}
\label{sec:ch15-curation}
% ----------------------------------------------------------------

\begin{theorem}[Curation Preserves Membership (15.1)]
\label{thm:curate-subset}
Every element of the curated repertoire was in the original:
\[
  \forall\, x \in \operatorname{curate}(\mathbf{R}, p),\quad x \in \mathbf{R}.
\]
\end{theorem}

\begin{proof}[Proof sketch]
By the property of list filtering: every element of
$\operatorname{filter}(p, \mathbf{R})$ belongs to $\mathbf{R}$.
\end{proof}

A museum can only display what it owns.  Curation is selection, not
creation---the formal basis of the archivist's constraint.  You cannot
curate what you do not have.  This seemingly obvious fact has deep
implications for cultural policy: any canon-formation exercise is bounded
by the existing cultural stock.  A national curriculum can only draw from
the ideas already present in the culture; it cannot conjure new ones through
the act of selection.\footnote{This distinguishes curation from
\emph{creation}.  In our model, creation is the introduction of a genuinely
new element $a \in \IS$ not already in the repertoire; curation is the
selection among existing elements.  The two operations have fundamentally
different structural properties.}

\begin{theorem}[Curation Reduces Length (15.2)]
\label{thm:curate-length}
The curated repertoire is never longer than the original:
\[
  |\operatorname{curate}(\mathbf{R}, p)| \;\le\; |\mathbf{R}|.
\]
\end{theorem}

\begin{proof}[Proof sketch]
Filtering can only remove elements, never add them.
\end{proof}

Curation only removes---it never inflates.  A censored library is always
smaller than the uncensored original.  This formalizes the ``information
loss'' inherent in any act of selection.  Every curatorial decision is,
structurally, a decision to \emph{forget}.

\begin{theorem}[Vacuous Curation is Empty (15.3)]
\label{thm:vacuous-curate}
Filtering by a predicate that rejects everything yields the empty list:
$\operatorname{curate}(\mathbf{R}, \lambda x.\, \texttt{false}) = []$.
\end{theorem}

\begin{proof}[Proof sketch]
No element passes the filter.
\end{proof}

Total censorship---rejecting every idea---produces cultural void.  This is
the formal limit of suppression: complete censorship equals complete cultural
erasure.  The totalitarian ideal of the \emph{tabula rasa}---erasing all
prior culture to build anew---is structurally equivalent to reducing the
repertoire to the empty list.  Pol Pot's ``Year Zero,'' Mao's Cultural
Revolution, the \emph{damnatio memoriae} of the Roman Senate: all are
attempts at vacuous curation.\footnote{The historical record suggests that
vacuous curation is never perfectly achieved: underground literature,
oral tradition, and material culture survive even the most thorough
purges.  The theorem describes the \emph{structural limit}, not the
empirical reality.}

\begin{theorem}[Universal Curation is Identity (15.4)]
\label{thm:universal-curate}
Filtering by a predicate that accepts everything yields the original:
$\operatorname{curate}(\mathbf{R}, \lambda x.\, \texttt{true}) = \mathbf{R}$.
\end{theorem}

\begin{proof}[Proof sketch]
Every element passes the filter.
\end{proof}

A museum that displays everything is just a warehouse.  ``Curating'' with no
selection criterion is structurally equivalent to no curation at all.  This
is the formal basis of the distinction between the \emph{archive}
(everything preserved) and the \emph{exhibition} (a curated selection).
The uncurated repertoire is the Borgesian library: complete, but
unnavigable.\footnote{Jorge Luis Borges, ``The Library of Babel'' (1941).
Borges's infinite library contains every possible book but is useless
precisely because it lacks curation.}

\begin{remark}[Connection to Information Theory: Channel Capacity]
\label{rem:channel-capacity}
The duality between curation (reducing the repertoire) and universal
acceptance (identity) has a precise analogue in Shannon's information
theory.  A communication channel with capacity $C$ can transmit at most
$C$ bits per second reliably; attempting to transmit more introduces
errors.  Similarly, a cognitive agent with finite processing capacity
can only maintain and effectively deploy a repertoire of bounded size.
The curatorial predicate $p$ functions as a \emph{channel code},
selecting which elements of the vast ideatic space $\IS$ to retain
within the agent's finite capacity.

Shannon's noisy channel coding theorem states that optimal codes exist
which achieve the capacity bound---but finding them is computationally
hard.  The analogous statement for curation is: for any cognitive
capacity $C$, there exists an optimal predicate $p^*$ that maximizes
$\operatorname{depthSum}(\operatorname{curate}(\mathbf{R}, p^*))$
subject to $|\operatorname{curate}(\mathbf{R}, p^*)| \le C$.  The
curator's dilemma is that this optimal predicate is, in general,
intractable to compute---hence the irreducibly judgmental character of
curatorial work.  No algorithm can replace the curator's taste, for
the same reason that no algorithm can efficiently find optimal channel
codes in the general case.
\end{remark}

\begin{example}[Spotify's Recommendation Algorithm as Automated Curation]
\label{ex:spotify}
Spotify's recommendation engine curates from a catalog of over 100
million tracks.  Its predicate $p_{\text{Spotify}}$ is a learned function
of the user's listening history, demographic data, audio features, and
collaborative filtering across millions of users.  In our framework,
each weekly ``Discover Weekly'' playlist is a curated repertoire:
$\mathbf{R}_{\text{playlist}} = \operatorname{curate}(
\mathbf{R}_{\text{catalog}}, p_{\text{Spotify}})$.

By Theorem~\ref{thm:curate-subset}, every recommended track was already
in the catalog; the algorithm creates nothing new.  By
Theorem~\ref{thm:curate-length}, the playlist is vastly shorter than the
catalog---typically 30 tracks from 100 million.  The crucial question is
whether the algorithmic predicate selects for depth or for familiarity.
Empirical evidence suggests that recommendation algorithms optimize for
engagement (clicks, completions, saves), which correlates imperfectly
with interpretive depth.  The result is a ``filter bubble'' in which the
user's curated repertoire converges to a narrow region of ideatic space,
reducing the diversity of interpretive possibilities available when
encountering novel signals.
\end{example}

% ----------------------------------------------------------------
\section{Expansion Fundamentals}
\label{sec:ch15-expansion}
% ----------------------------------------------------------------

\begin{theorem}[Expansion Increases Length (15.5)]
\label{thm:expand-length}
Adding a new idea increases the repertoire by exactly one:
\[
  |\operatorname{expand}(\mathbf{R}, a)| = |\mathbf{R}| + 1.
\]
\end{theorem}

\begin{proof}[Proof sketch]
Prepending one element increases list length by one.
\end{proof}

Each learning event adds exactly one new schema.  Education is
\emph{incremental}---you learn one thing at a time.  This constrains
theories of ``epiphany'' or ``gestalt shift'' to be single-idea additions
in the formal model.  The Humboldt-ian ideal of \emph{Bildung}---the
gradual, systematic expansion of the mind---is structurally captured by
iterated application of $\operatorname{expand}$.

\begin{theorem}[New Idea is in Expanded Repertoire (15.6)]
\label{thm:new-idea-mem}
The newly added idea belongs to the expanded repertoire:
$a \in \operatorname{expand}(\mathbf{R}, a)$.
\end{theorem}

\begin{proof}[Proof sketch]
The new idea is the head of the list.
\end{proof}

What you learn, you possess.  The student who has been taught a concept
\emph{has} that concept in their repertoire.  This is the structural basis
of education as repertoire expansion: successful pedagogy is precisely the
operation $\operatorname{expand}(\mathbf{R}_{\text{student}}, a_{\text{lesson}})$.

\begin{theorem}[Old Ideas Survive Expansion (15.7)]
\label{thm:old-survive}
All ideas in the original repertoire remain in the expanded version:
\[
  \forall\, x \in \mathbf{R},\quad x \in \operatorname{expand}(\mathbf{R}, a).
\]
\end{theorem}

\begin{proof}[Proof sketch]
The original list is the tail of the expanded list.
\end{proof}

Learning does not erase prior knowledge.  The student who learns calculus
still remembers arithmetic.  Expansion is \emph{monotone}---repertoires only
grow.  This is an important modeling assumption: in our framework, forgetting
requires an explicit curation step (filtering out the unwanted idea), not
mere neglect.  The distinction between passive forgetting and active
censorship is, in IDT, the distinction between non-operation and
curation.\footnote{In cognitive science, the debate between ``decay'' and
``interference'' theories of forgetting maps onto our distinction between
curation (active removal) and repertoire stagnation (non-expansion).}

% ----------------------------------------------------------------
\section{Curation and Interpretation}
\label{sec:ch15-curate-interp}
% ----------------------------------------------------------------

\begin{theorem}[Curated Interpretation is Sublist (15.8)]
\label{thm:curated-interp}
Interpreting a curated repertoire gives at most as many interpretations as
the full repertoire:
\[
  |\operatorname{interpret}(\operatorname{curate}(\mathbf{R}, p), s)|
  \;\le\;
  |\operatorname{interpret}(\mathbf{R}, s)|.
\]
\end{theorem}

\begin{proof}[Proof sketch]
$|\operatorname{interpret}(\cdot, s)| = |\cdot|$ (interpretation preserves
length), and $|\operatorname{curate}(\mathbf{R}, p)| \le |\mathbf{R}|$.
\end{proof}

A censored mind produces fewer interpretations.  The number of possible
readings of a text is bounded by the reader's (curated) repertoire size.
A narrower education yields fewer perspectives---the formal basis of the
humanist argument for breadth in the curriculum.  The liberally educated
mind, with its expansive repertoire, generates more interpretations than
the narrowly trained specialist.

\begin{remark}[Censorship as Curation with Political Consequences]
\label{rem:censorship-political}
The connection between curation and censorship deserves sustained
attention, because the formal equivalence conceals a crucial normative
distinction.  The museum curator and the government censor both apply
predicates to repertoires, but they differ in three respects: (1) the
scope of application (one collection vs.\ an entire society), (2) the
coercive power backing the predicate (aesthetic judgment vs.\ legal
penalty), and (3) the reversibility of the operation (the deaccessioned
painting can be reacquired; the burned book is gone forever).

Our formal framework captures the \emph{structural} identity of these
operations while remaining silent on their normative valence.  This is
a feature, not a bug: the theorems show that the same mathematical
constraints govern benign curation and malign censorship, which means
that arguments about the \emph{structural} effects of censorship need
not appeal to contested moral premises.
Theorem~\ref{thm:curated-interp}---that curation reduces interpretive
capacity---is value-neutral as a mathematical fact, but devastating as
an argument against censorship: it shows that \emph{any} censorship,
regardless of its motivation, reduces the society's collective capacity
to interpret signals.
\end{remark}

\begin{theorem}[Empty Curation Yields No Interpretation (15.9)]
\label{thm:empty-curate-interp}
If curation removes everything, there are no interpretations:
\[
  \operatorname{interpret}(\operatorname{curate}(\mathbf{R},
  \lambda x.\, \texttt{false}), s) = [].
\]
\end{theorem}

\begin{proof}[Proof sketch]
By Theorem~\ref{thm:vacuous-curate}, the curated repertoire is empty,
and mapping over the empty list yields the empty list.
\end{proof}

The culturally emptied mind cannot interpret anything.  Total censorship
followed by signal reception equals zero interpretive output.  This is the
``blank slate after brainwashing'' theorem: the totalitarian regime that
succeeds in erasing all prior culture leaves its subjects incapable of
interpreting \emph{any} signal, including its own propaganda.  The self-
defeating nature of complete censorship is a structural theorem.

% ----------------------------------------------------------------
\section{Expansion and Interpretation}
\label{sec:ch15-expand-interp}
% ----------------------------------------------------------------

\begin{theorem}[Expanded Interpretation Adds One (15.10)]
\label{thm:expanded-interp}
Interpreting an expanded repertoire produces one more interpretation than
the original:
\[
  |\operatorname{interpret}(\operatorname{expand}(\mathbf{R}, a), s)|
  \;=\;
  |\operatorname{interpret}(\mathbf{R}, s)| + 1.
\]
\end{theorem}

\begin{proof}[Proof sketch]
Expansion adds one element; interpretation preserves length.
\end{proof}

Each new idea in the repertoire adds exactly one new interpretive
possibility.  Education's return is precisely one new perspective per concept
learned---the fundamental unit of \emph{Bildung}.  The marginal return on
education, measured in interpretive possibilities, is constant: one concept
in, one new interpretation out.

\begin{theorem}[Expansion Interpretation Structure (15.11)]
\label{thm:expand-structure}
Interpreting an expanded repertoire prepends the new interpretation:
\[
  \operatorname{interpret}(\operatorname{expand}(\mathbf{R}, a), s)
  \;=\;
  (a \compose s) :: \operatorname{interpret}(\mathbf{R}, s).
\]
\end{theorem}

\begin{proof}[Proof sketch]
By definition: $\operatorname{expand}(\mathbf{R}, a) = a :: \mathbf{R}$,
and mapping over a cons is $f(a) :: \operatorname{map}(f, \mathbf{R})$.
\end{proof}

The new concept's interpretation appears first, followed by all prior
interpretations unchanged.  Learning adds a new lens without distorting
existing ones---the ``additive'' model of education.  This is structurally
optimistic: new knowledge does not corrupt old knowledge.  The Baconian
vision of knowledge as cumulative finds formal support here.

\begin{theorem}[Void Expansion Doesn't Change Interpretations (15.12)]
\label{thm:void-expansion}
Adding void to the repertoire adds the raw signal $s$ as a new interpretation:
\[
  \operatorname{interpret}(\operatorname{expand}(\mathbf{R}, \void), s)
  \;=\;
  s :: \operatorname{interpret}(\mathbf{R}, s).
\]
\end{theorem}

\begin{proof}[Proof sketch]
$\void \compose s = s$ by the void axiom, so the new interpretation is $s$
itself.
\end{proof}

Adding ``empty receptivity''---the Keatsian \emph{negative capability}---to
one's repertoire adds the raw signal as a new interpretation.  The
educationally valuable ``empty slot'' lets you see the signal unmediated.
This is the formal justification for intellectual humility: maintaining a
void element in one's repertoire ensures that one always has access to the
unfiltered signal, free from the bias of pre-existing categories.\footnote{Keats's
concept of ``negative capability''---``being in uncertainties, mysteries,
doubts, without any irritable reaching after fact and reason''---is
precisely the act of maintaining $\void$ in the repertoire: a position of
receptivity that produces the raw signal as interpretation.}

\begin{lstlisting}[caption={Void expansion raw signal in Lean~4}]
theorem void_expansion_raw_signal (rep : List I) (s : I) :
    interpret (expand rep void) s = s :: interpret rep s := by
  simp [expanded_interp_structure, void_left]
\end{lstlisting}

% ----------------------------------------------------------------
\section{DepthSum Bounds}
\label{sec:ch15-depthsum}
% ----------------------------------------------------------------

\begin{theorem}[Curated DepthSum Bound (15.13)]
\label{thm:curated-depthsum}
The depth sum of a curated repertoire is at most the depth sum of the
original:
\[
  \operatorname{depthSum}(\operatorname{curate}(\mathbf{R}, p))
  \;\le\;
  \operatorname{depthSum}(\mathbf{R}).
\]
\end{theorem}

\begin{proof}[Proof sketch]
By induction on $\mathbf{R}$.  At each step, if the element passes the
filter, both sides include $\depth(r)$; if it fails, only the right side
includes it, maintaining the inequality.
\end{proof}

Curation can only \emph{reduce} total cultural depth.  A museum that hides
artifacts loses depth; a censored library is shallower than the uncensored
one.  Information destruction is irreversible in the depth metric.  This
theorem has stark implications for cultural policy: every act of censorship
is a permanent reduction in the society's total interpretive capacity.  The
censored ideas' depth is \emph{lost}, not merely hidden.

\begin{theorem}[Expanded DepthSum (15.14)]
\label{thm:expanded-depthsum}
Expanding the repertoire adds exactly the new idea's depth:
\[
  \operatorname{depthSum}(\operatorname{expand}(\mathbf{R}, a))
  \;=\;
  \depth(a) + \operatorname{depthSum}(\mathbf{R}).
\]
\end{theorem}

\begin{proof}[Proof sketch]
$\operatorname{expand}(\mathbf{R}, a) = a :: \mathbf{R}$, and
$\operatorname{depthSum}(a :: \mathbf{R}) = \depth(a) +
\operatorname{depthSum}(\mathbf{R})$.
\end{proof}

Each new idea adds exactly its depth to the total cultural capital.
Education is a \emph{linear} accumulation of depth---Bourdieu's cultural
capital grows additively with each new acquisition.  The student who learns
a concept of depth $d$ increases their total cultural capital by exactly $d$.
This linearity is both empowering (every lesson counts) and sobering (there
are no shortcuts---depth must be earned one idea at a time).

\begin{theorem}[Void Expansion DepthSum (15.15)]
\label{thm:void-expansion-depth}
Adding void to the repertoire does not change depthSum:
\[
  \operatorname{depthSum}(\operatorname{expand}(\mathbf{R}, \void))
  \;=\;
  \operatorname{depthSum}(\mathbf{R}).
\]
\end{theorem}

\begin{proof}[Proof sketch]
$\depth(\void) = 0$, so $0 + \operatorname{depthSum}(\mathbf{R}) =
\operatorname{depthSum}(\mathbf{R})$.
\end{proof}

Adding ``nothing'' to one's cultural stock adds zero depth.  The void is
costless to maintain in the repertoire---the formal basis of ``keeping an
open mind'' being free in the depth metric.  Intellectual humility (carrying
$\void$ in one's repertoire) costs nothing and gains the benefit of raw
signal access (Theorem~\ref{thm:void-expansion}).

\begin{example}[Museum Deaccessioning Controversies]
\label{ex:deaccessioning}
Museum deaccessioning---the practice of removing works from a permanent
collection, typically through sale---illustrates the irreversibility of
curation in practice.  When the Berkshire Museum in 2017 proposed selling
Norman Rockwell paintings and Hudson River School landscapes to fund
operational costs, the art world erupted.  The controversy turns on a
formal point: deaccessioning applies a new predicate
$p_{\text{financial}}$ (``does this work justify its storage and
insurance costs?'') to a repertoire that was originally curated by a
different predicate $p_{\text{aesthetic}}$ (``does this work have
artistic merit?'').

By Theorem~\ref{thm:curate-idempotent}, re-curating by the same
predicate has no effect.  But curating by a \emph{different} predicate
compounds the restriction:
$\operatorname{curate}(\operatorname{curate}(\mathbf{R},
p_{\text{aesthetic}}), p_{\text{financial}})$ yields a repertoire that
satisfies \emph{both} predicates, and is therefore (in general) strictly
smaller than either single-predicate curation.  The lost works satisfy
$p_{\text{aesthetic}}$ but not $p_{\text{financial}}$; they are
aesthetically valuable but financially burdensome.  The museum's
cultural depth (Theorem~\ref{thm:curated-depthsum}) shrinks by the sum
of the deaccessioned works' depths.

The deaccessioning debate thus reduces to a question about predicate
ordering: should financial predicates ever override aesthetic ones in
cultural institutions?  Our framework does not answer this normative
question, but it precisely quantifies the cost: each deaccessioned work
reduces $\operatorname{depthSum}$ by exactly $\depth(r_{\text{sold}})$,
and this reduction is permanent once the work enters a private collection.
\end{example}

% ----------------------------------------------------------------
\section{Structural Curation Theorems}
\label{sec:ch15-structural}
% ----------------------------------------------------------------

\begin{theorem}[Double Curation is Single Curation (15.16)]
\label{thm:curate-idempotent}
Curating a curated repertoire by the same predicate is idempotent:
\[
  \operatorname{curate}(\operatorname{curate}(\mathbf{R}, p), p)
  \;=\;
  \operatorname{curate}(\mathbf{R}, p).
\]
\end{theorem}

\begin{proof}[Proof sketch]
$\operatorname{filter}(p, \operatorname{filter}(p, \mathbf{R}))
= \operatorname{filter}(p \wedge p, \mathbf{R})
= \operatorname{filter}(p, \mathbf{R})$
since $p \wedge p = p$.
\end{proof}

Censoring an already-censored library by the same criterion has no further
effect.  The ``diminishing returns'' of redundant censorship---once you have
removed all ``heretical'' books, removing them again changes nothing.  This
idempotency result has practical implications for bureaucratic censorship:
the second censor applying the same criteria as the first is structurally
redundant, a waste of institutional resources.  Totalitarian regimes that
layer censorship committees atop each other (with identical mandates) achieve
nothing beyond the first pass.\footnote{Historical examples include the
multiple overlapping censorship bodies of the Soviet Union (\textit{Glavlit},
\textit{Goskomizdat}, Party committees), each applying essentially the same
ideological filter.  Our theorem shows this redundancy is structural, not
merely bureaucratic.}

\begin{theorem}[Expand Then Curate Accepting (15.17)]
\label{thm:expand-curate-accept}
If the new idea passes the predicate, expanding then curating keeps it:
\[
  p(a) = \texttt{true}
  \;\Longrightarrow\;
  \operatorname{curate}(\operatorname{expand}(\mathbf{R}, a), p)
  = a :: \operatorname{curate}(\mathbf{R}, p).
\]
\end{theorem}

\begin{proof}[Proof sketch]
Filtering $a :: \mathbf{R}$ with $p$: since $p(a) = \texttt{true}$,
$a$ is kept, and the tail is filtered as usual.
\end{proof}

If a newly acquired idea passes the censorship test, it survives curation.
``Approved learning'' is permanent---the accepted idea stays in the
repertoire through subsequent curation rounds.  This is the formal basis
of canonical inclusion: once an idea is admitted to the canon (passes the
predicate), it persists through all future curation by the same criterion.
Dante, admitted to the Western canon, survives every subsequent round of
canon revision that uses the same evaluative framework.

\begin{theorem}[Expand Then Curate Rejecting (15.18)]
\label{thm:expand-curate-reject}
If the new idea fails the predicate, expanding then curating removes it:
\[
  p(a) = \texttt{false}
  \;\Longrightarrow\;
  \operatorname{curate}(\operatorname{expand}(\mathbf{R}, a), p)
  = \operatorname{curate}(\mathbf{R}, p).
\]
\end{theorem}

\begin{proof}[Proof sketch]
Filtering $a :: \mathbf{R}$ with $p$: since $p(a) = \texttt{false}$,
$a$ is discarded, and the tail is filtered as before.
\end{proof}

Ideas that fail the censorship test are immediately purged.  The
``forbidden knowledge'' theorem: learning something that the regime rejects
is structurally equivalent to never having learned it, after curation.  This
formalizes the futility of teaching banned concepts in totalitarian
educational systems---the student may \emph{encounter} the forbidden idea
(via $\operatorname{expand}$), but the subsequent curation step erases it
from the repertoire as if the encounter never happened.

The theorem is disturbing in its implications.  If $p$ represents the
ideological filter of a totalitarian state, then
$\operatorname{curate}(\operatorname{expand}(\mathbf{R}, a_{\text{forbidden}}), p)
= \operatorname{curate}(\mathbf{R}, p)$: the forbidden idea leaves no trace.
The dissident's effort to teach subversive ideas is, under perfect
censorship, structurally null.  Of course, perfect censorship is an
idealization---but the theorem shows what perfect censorship would
\emph{mean}: structural erasure of dissent.

\paragraph{Extended proof discussion.}
The expand-then-curate-rejecting theorem (15.18) merits careful
reflection on its proof structure, because it reveals a deep asymmetry
between expansion and curation.  The proof proceeds by case analysis on
the head of the list $a :: \mathbf{R}$.  When $p(a) = \texttt{false}$,
the filter function skips the head and recurses on the tail, producing
exactly $\operatorname{filter}(p, \mathbf{R})$.  The key insight is that
the filter is \emph{memoryless}: it does not record that an element was
tested and rejected.  There is no ``trace'' of the rejected idea in the
output.  This memorylessness is what makes censorship structurally
complete in the formal model: the curated repertoire contains no evidence
that anything was removed.

Contrast this with real-world censorship, where the act of banning a
book often draws attention to it (the ``Streisand effect'').  The
discrepancy arises because real censorship operates in a social context
where the \emph{act} of curation is itself observable as a signal.  In
our model, the curatorial operation is transparent to downstream
consumers of the repertoire: they see only
$\operatorname{curate}(\mathbf{R}, p)$, not the pair
$(\mathbf{R}, p)$.  Modeling censorship-awareness would require a
second-order theory in which agents can observe and interpret curatorial
predicates as signals---a direction we leave for future work.

\begin{remark}[The Filter Bubble Problem]
\label{rem:filter-bubble}
Eli Pariser's concept of the ``filter bubble'' (\textit{The Filter
Bubble}, 2011) describes the phenomenon in which algorithmic curation
progressively narrows an individual's information environment.  In our
formal terms, a filter bubble arises when the curatorial predicate is
\emph{adaptive}: $p_{t+1}$ is a function of
$\operatorname{curate}(\mathbf{R}, p_t)$.  If the algorithm curates
based on what the user has previously consumed, and the user consumes
only what is curated, a feedback loop emerges.

Formally, let $p_0$ be the initial predicate, and define
$\mathbf{R}_{t+1} = \operatorname{curate}(\mathbf{R}_t, p_t)$ where
$p_t$ depends on $\mathbf{R}_{t-1}$.  By Theorem~\ref{thm:curate-length},
$|\mathbf{R}_{t+1}| \le |\mathbf{R}_t|$, so the sequence
$\{|\mathbf{R}_t|\}$ is non-increasing.  If $p_t$ ever strictly
rejects an element of $\mathbf{R}_t$, the repertoire shrinks
\emph{irreversibly} (assuming no external expansion).  The filter
bubble is thus a monotone descent in repertoire size: once an idea is
filtered out, the adaptive predicate---which learns from the filtered
repertoire---has no way to ``rediscover'' it.  This is the formal
diagnosis of the filter bubble: adaptive curation without external
expansion converges to a fixed point of minimal repertoire that merely
confirms the user's existing ideas.

Pariser's solution---deliberate exposure to diverse viewpoints---is
formally an $\operatorname{expand}$ operation that injects ideas outside
the current curated repertoire, breaking the feedback loop by
introducing elements that the adaptive predicate has not yet evaluated.
The interplay between Theorems~\ref{thm:expand-curate-accept}
and~\ref{thm:expand-curate-reject} then determines whether these
injected ideas survive subsequent curation rounds.
\end{remark}

\begin{lstlisting}[caption={Expand-then-curate in Lean~4}]
theorem expand_curate_removes {new_idea : I} {p : I -> Bool}
    (hp : p new_idea = false) (rep : List I) :
    curate (expand rep new_idea) p = curate rep p := by
  simp [expand, curate, List.filter, hp]
\end{lstlisting}

% ----------------------------------------------------------------
\section{The Politics of Repertoire}
\label{sec:ch15-politics}
% ----------------------------------------------------------------

The theorems of this chapter, though stated in the neutral language of
predicates and lists, have inescapably political implications.  Every act
of curation is an exercise of power: the power to determine what ideas
are available for interpretation, and therefore what interpretations are
possible.  This section draws out the political consequences of the
formal framework.

\subsection{Who Holds the Predicate?}

The central political question of repertoire curation is: \emph{who
chooses the predicate $p$?}  In a democratic society, the answer is
ideally ``the community of interpreters themselves,'' through mechanisms
such as democratic curriculum design, public library governance, and
editorial independence.  In authoritarian regimes, the predicate is
imposed by the state: a single party, a censorship bureau, or a
supreme leader determines what ideas are admitted to the national
repertoire.

The formal framework reveals that this difference is not merely
quantitative (how much censorship) but qualitative (who controls the
structural operation).  The predicate $p$ is a \emph{function}, and the
choice of function space matters: a community that can choose from a
rich space of predicates (diverse curricula, multiple publishers,
independent media) has a structurally different curatorial capacity than
one restricted to a single state-imposed predicate.

\subsection{Curation and Inequality}

Bourdieu's analysis of cultural capital implies that curation is
unevenly distributed.  The children of the bourgeoisie enter school
with repertoires already curated by family exposure to art, literature,
and cosmopolitan experience.  The children of the working class enter
with repertoires curated by material necessity.  The school then applies
a uniform predicate---the curriculum---to both groups, but the
\emph{expansion} operations are different: the bourgeois child
encounters new ideas that resonate with an already-curated repertoire,
while the working-class child encounters ideas for which no resonating
elements exist in their current repertoire.

Formally, if $\mathbf{R}_{\text{bourgeois}}$ already contains elements
that resonate with the curriculum's new ideas, then expansion is
``sticky''---the new ideas compose meaningfully with existing ones,
producing high-depth interpretations.  If
$\mathbf{R}_{\text{working}}$ lacks such resonating elements, expansion
still adds the new idea (Theorem~\ref{thm:new-idea-mem}), but the
interpretive compositions it enables may have lower depth due to the
absence of resonant partners.  Educational inequality is thus a
\emph{compositional} phenomenon: it arises not from differential access
to ideas per se, but from differential resonance between new ideas and
existing repertoires.

\subsection{Repertoire Diversity as Public Good}

A society's collective repertoire---the union of all individual
repertoires---determines its collective interpretive capacity.  If every
citizen curates by the same predicate (ideological conformity), the
collective repertoire is no larger than a single individual's.  If
citizens curate by diverse predicates (ideological pluralism), the
collective repertoire is the union of many differently curated subsets,
and is therefore larger.

This suggests a formal argument for intellectual diversity as a public
good: a society with diverse curatorial predicates has a strictly larger
collective repertoire, and therefore a strictly greater collective
capacity to interpret novel signals.  When an unprecedented crisis
arrives---a pandemic, an economic shock, a technological disruption---the
society with a diverse repertoire has more interpretive resources to draw
upon than the homogeneous society.  Censorship, which imposes a uniform
predicate, is thus a structural risk factor for societal fragility:
it reduces the diversity of interpretive resources available when they
are most needed.

\paragraph{Extended proof discussion: Depth bounds and curatorial
composition.}
The interaction between depth bounds and iterated curation
deserves more detailed analysis.  Consider a sequence of curations
$\operatorname{curate}(\operatorname{curate}(\mathbf{R}, p_1), p_2)$.
By Theorem~\ref{thm:curated-depthsum} applied twice:
\[
  \operatorname{depthSum}(\operatorname{curate}(
    \operatorname{curate}(\mathbf{R}, p_1), p_2))
  \;\le\;
  \operatorname{depthSum}(\operatorname{curate}(\mathbf{R}, p_1))
  \;\le\;
  \operatorname{depthSum}(\mathbf{R}).
\]
Each curation pass can only decrease depth, so iterated curation by
different predicates produces a monotone-decreasing sequence of depth
sums.  This means that the order of application does not matter for the
\emph{direction} of the inequality (both orderings lose depth), but it
may matter for the \emph{amount} of depth lost.

Specifically, if $p_1$ removes deep elements and $p_2$ removes shallow
elements, the total depth loss is the sum of the depths of all removed
elements.  But if $p_1$ removes some elements that $p_2$ would also have
removed, the compound curation $\operatorname{curate}(\cdot, p_1 \wedge
p_2)$---applying both predicates simultaneously---may yield a different
depth sum than either sequential application.  The key structural
insight is that $\operatorname{curate}(\operatorname{curate}(\mathbf{R},
p_1), p_2) = \operatorname{curate}(\mathbf{R}, p_1 \wedge p_2)$
when both predicates are applied to the same elements, which follows
from the fact that $\operatorname{filter}(p_2,
\operatorname{filter}(p_1, \mathbf{R})) =
\operatorname{filter}(p_1 \wedge p_2, \mathbf{R})$ for any list
$\mathbf{R}$.  This filter-conjunction lemma is the formal basis of the
observation that layered censorship (first political, then moral, then
aesthetic) is equivalent to a single pass with a conjunctive predicate.

% ----------------------------------------------------------------
\section{Conclusion}
\label{sec:ch15-conclusion}
% ----------------------------------------------------------------

Repertoire curation, formalized in IDT, reveals the structural anatomy of
cultural management.  Curation is subtractive
(Theorem~\ref{thm:curate-length}), idempotent
(Theorem~\ref{thm:curate-idempotent}), and depth-reducing
(Theorem~\ref{thm:curated-depthsum}).  Expansion is additive
(Theorem~\ref{thm:expand-length}), monotone
(Theorem~\ref{thm:old-survive}), and depth-increasing
(Theorem~\ref{thm:expanded-depthsum}).  The interaction of the two---expand
then curate---is determined entirely by whether the new idea passes the
filter (Theorems~\ref{thm:expand-curate-accept}
and~\ref{thm:expand-curate-reject}).

These results connect to four domains of cultural practice.  In
\emph{museum theory}, curation is the fundamental operation of exhibition
design---and our theorems show that it is structurally subtractive.  In
\emph{education policy}, expansion is the mechanism of \emph{Bildung},
and each lesson adds exactly one interpretive possibility.  In
\emph{censorship dynamics}, vacuous curation produces cultural void
(Theorem~\ref{thm:vacuous-curate}), and the expand-then-curate-rejecting
result (Theorem~\ref{thm:expand-curate-reject}) formalizes the structural
erasure of dissent.  In \emph{Bourdieu's sociology}, cultural capital grows
linearly with expansion (Theorem~\ref{thm:expanded-depthsum}) and can only
shrink through curation (Theorem~\ref{thm:curated-depthsum}).

The five chapters of this sequence---marketplace competition
(Chapter~\ref{ch:marketplace}), stable matching
(Chapter~\ref{ch:matching}), coalition value
(Chapter~\ref{ch:coalition}), ideatic auction
(Chapter~\ref{ch:auction}), and repertoire curation---together provide a
comprehensive game-theoretic framework for cultural dynamics within
Ideatic Diffusion Theory.  Each chapter isolates a different mode of
cultural interaction---competition, pairing, cooperation, bidding,
management---and derives structural constraints that hold across all
instantiations of the IDT axioms.  The resulting picture is one of
\emph{bounded cultural possibility}: every mode of interaction is
constrained by subadditivity, identity, and resonance, the three pillars
on which the theory rests.
